\documentclass[12pt,a4paper]{article}
\usepackage[a4paper,margin=18mm]{geometry}
\usepackage[T2A]{fontenc}
\usepackage[utf8]{inputenc}
\usepackage[russian]{babel}
\usepackage{amsmath,amssymb,mathtools}
\usepackage{array,booktabs,longtable}
\usepackage{xcolor}
\usepackage{hyperref}
\usepackage{tikz}

\hypersetup{
  colorlinks=true,
  linkcolor=blue!55!black,
  urlcolor=blue!55!black
}

\setlength{\parindent}{0pt}
\setlength{\parskip}{6pt}
\setlength{\tabcolsep}{3pt}
\renewcommand{\arraystretch}{1.18}
\sloppy

\begin{document}

\pagenumbering{roman}
\begin{titlepage}
\centering
\vspace*{0.27\textheight}

{\LARGE\bfseries Ревью домашней работы\par}

\vspace{1.6cm}

{\large Автор: Лёвин Артём Александрович\par}

\vspace{0.8cm}

{\large 14 сентября 2026 года\par}

\vfill

\begin{tikzpicture}
  \draw[blue!45!black, line width=0.8pt]
    (0,0) .. controls (1.5,0.45) and (3.5,-0.45) .. (5,0)
          .. controls (6.5,0.45) and (8.5,-0.45) .. (10,0);
\end{tikzpicture}

\end{titlepage}

\pagenumbering{arabic}
\newpage
\section*{Сводная таблица}
\small

\begin{tabular}{@{}
>{\raggedright\arraybackslash}p{0.08\linewidth}
>{\raggedright\arraybackslash}p{0.16\linewidth}
>{\raggedright\arraybackslash}p{0.29\linewidth}
>{\raggedright\arraybackslash}p{0.17\linewidth}
>{\raggedright\arraybackslash}p{0.18\linewidth}
@{}}
\toprule
№ задания & Тема & Суть ошибки & Ваш ответ & Правильный ответ \\
\midrule
\hyperref[task:2]{2} &
Степень степени и произведение степеней &
Сначала показатели степеней сложены вместо умножения, затем показатели перемножены вместо сложения. Две ошибки дали правильный финальный показатель. &
$x^5$ &
$x^5$, $x\ne0$ \\
\addlinespace
\hyperref[task:4]{4} &
Квадратный корень из квадрата &
При извлечении корня потерян модуль, поэтому при отрицательном $x$ получен неверный знак; дополнительно неверно вычислено $(-2)^5$. &
последняя запись $81\cdot(-512)$ &
$2592$ \\
\addlinespace
\hyperref[task:5]{5} &
Свойства степеней при делении &
После корректного раскрытия квадрата и сокращения множителей оставшаяся степень $a^5$ потеряна; записан ноль. &
$0$ &
$a^5$, $a\ne0$, $b\ne0$ \\
\addlinespace
\hyperref[task:6]{6} &
Квадратный корень из полного квадрата &
Записано $\sqrt{(6x-5)^2}=6x-5$ вместо модуля. При заданном $x$ выражение внутри модуля отрицательно. &
числовой ответ не указан &
$2$ \\
\addlinespace
\hyperref[task:8]{8} &
Отрицательные показатели степени &
Показатели при делении степеней обработаны верно, но $(-2)^5$ ошибочно вычислено как $-512$. &
$-512$ &
$-32$ \\
\addlinespace
\hyperref[task:9]{9} &
Радикалы и модуль &
Использовано $\sqrt{b^6}=b^3$ при $b=-2$. Должно быть $|b^3|$; две одинаковые ошибки в числителе и знаменателе взаимно компенсировались. &
$300$ &
$300$ \\
\bottomrule
\end{tabular}

\normalsize

\newpage
\null
\vfill
\begin{center}
{\LARGE\bfseries Разбор ошибок}
\end{center}
\vfill
\null

\newpage
\phantomsection
\label{task:2}
\section*{Задание №2}

\textbf{Текст задания.}

Упростите выражение
\[
\frac{(x^4)^3x^2}{x^9}.
\]

\textbf{Ваше решение.}

В работе записана цепочка
\[
\frac{(x^4)^3x^2}{x^9}
=
\frac{x^7x^2}{x^9}
=
\frac{x^{14}}{x^9}
=
x^5.
\]

\textbf{Ваш ответ.}

\[
x^5.
\]

\textcolor{red}{Ошибка:} перепутаны два разных правила действий со степенями.

\textbf{Где именно возникает ошибка.}

Последний корректный шаг --- исходная запись
\[
\frac{(x^4)^3x^2}{x^9}.
\]

Первый неверный переход:
\[
(x^4)^3=x^7.
\]
Здесь показатели $4$ и $3$ были сложены. Для степени степени это правило не подходит.

Следом появляется вторая ошибка:
\[
x^7\cdot x^2=x^{14}.
\]
Здесь показатели, наоборот, были перемножены, хотя при умножении степеней с одинаковым основанием показатели нужно складывать.

\textbf{Почему это неверно.}

Есть два разных правила, и их важно не менять местами.

Для степени степени:
\[
(x^m)^n=x^{mn}.
\]
Поэтому
\[
(x^4)^3=x^{4\cdot3}=x^{12}.
\]

Для произведения степеней с одинаковым основанием:
\[
x^m\cdot x^n=x^{m+n}.
\]
Поэтому
\[
x^{12}\cdot x^2=x^{12+2}=x^{14}.
\]

В работе две разные ошибки случайно привели к показателю $14$, который действительно должен появиться в числителе. Такой совпавший промежуточный результат не делает ход решения правильным.

Кроме того, исходный знаменатель содержит $x^9$, поэтому необходимо учитывать условие
\[
x\ne0.
\]

\textcolor{blue!60!black}{Ключевая идея:} внутри конструкции ``степень степени'' показатели перемножаются, а при умножении одинаковых оснований показатели складываются.

\textbf{Исправление от последнего правильного шага.}

Начинаем с исходного выражения и применяем правила в правильном порядке:
\[
\frac{(x^4)^3x^2}{x^9}
=
\frac{x^{12}x^2}{x^9}
=
\frac{x^{14}}{x^9}
=
x^{14-9}
=
x^5.
\]

\textbf{Полное правильное решение.}

Так как в знаменателе находится $x^9$, имеем $x\ne0$.

Сначала возводим степень в степень:
\[
(x^4)^3=x^{4\cdot3}=x^{12}.
\]
Далее умножаем степени с одинаковым основанием:
\[
x^{12}x^2=x^{12+2}=x^{14}.
\]
Теперь делим степени с одинаковым основанием:
\[
\frac{x^{14}}{x^9}=x^{14-9}=x^5.
\]

\textbf{Проверка.}

Возьмём, например, допустимое значение $x=2$. Исходное выражение равно
\[
\frac{(2^4)^3\cdot2^2}{2^9}
=
\frac{2^{12}\cdot2^2}{2^9}
=
2^5.
\]
Это совпадает с полученным результатом.

\textcolor{teal!60!black}{Итог:}
\[
\boxed{x^5},\qquad x\ne0.
\]

\textbf{Задача на закрепление.}

Упростите выражение
\[
\frac{(m^5)^2m^3}{m^6},\qquad m\ne0.
\]

\newpage
\phantomsection
\label{task:4}
\section*{Задание №4}

\textbf{Текст задания.}

Найдите значение
\[
\sqrt{81x^{10}y^4}
\]
при $x=-2$, $y=3$.

\textbf{Ваше решение.}

В работе записано
\[
\sqrt{81x^{10}y^4}
=
9x^5y^2
=
9(-2)^5(3)^2,
\]
после чего появляется вычисление
\[
9\cdot(-512)\cdot9=81\cdot(-512).
\]

\textbf{Ваш ответ.}

Окончательный числовой ответ не доведён; последняя запись --- $81\cdot(-512)$.

\textcolor{red}{Ошибка:} при извлечении квадратного корня из квадрата потерян модуль; затем неверно вычислена пятая степень числа $-2$.

\textbf{Где именно возникает ошибка.}

Последний корректный шаг:
\[
\sqrt{81x^{10}y^4}.
\]

Первый неверный переход:
\[
\sqrt{81x^{10}y^4}=9x^5y^2.
\]
Поскольку
\[
x^{10}=(x^5)^2,
\]
из квадратного корня должно появиться не $x^5$, а $|x^5|$.

При заданном $x=-2$ это существенно, потому что
\[
x^5=(-2)^5<0.
\]
Квадратный корень по определению не может быть отрицательным.

Отдельно в работе записано
\[
(-2)^5=-512,
\]
но это также неверно: число $512$ получается как $2^9$, а не как $2^5$.

\textbf{Почему это неверно.}

Основное правило:
\[
\sqrt{u^2}=|u|,
\]
а не просто $u$. Модуль необходим потому, что квадратный корень обозначает неотрицательное число.

Здесь
\[
\sqrt{x^{10}}=\sqrt{(x^5)^2}=|x^5|.
\]
Кроме того,
\[
\sqrt{y^4}=\sqrt{(y^2)^2}=|y^2|=y^2,
\]
поскольку $y^2\ge0$ при любом действительном $y$.

Поэтому
\[
\sqrt{81x^{10}y^4}=9|x^5|y^2.
\]

При $x=-2$:
\[
(-2)^5=-32,
\qquad
|(-2)^5|=32.
\]

\textcolor{blue!60!black}{Ключевая идея:} если из квадратного корня извлекается квадрат выражения, нужно проверить его знак и в общем случае поставить модуль.

\textbf{Исправление от последнего правильного шага.}

\[
\sqrt{81x^{10}y^4}
=
9|x^5|y^2.
\]
Подставляем $x=-2$, $y=3$:
\[
9|(-2)^5|\cdot3^2
=
9\cdot32\cdot9
=
2592.
\]

\textbf{Полное правильное решение.}

Представим подкоренное выражение как произведение квадратов:
\[
81x^{10}y^4
=
9^2(x^5)^2(y^2)^2.
\]
Тогда
\[
\sqrt{81x^{10}y^4}
=
\sqrt{9^2}\,\sqrt{(x^5)^2}\,\sqrt{(y^2)^2}
=
9|x^5|y^2.
\]
Подставляем данные значения:
\[
9|(-2)^5|\cdot3^2
=
9\cdot32\cdot9
=
2592.
\]

\textbf{Проверка.}

Можно подставить значения сразу под корень:
\[
81(-2)^{10}3^4
=
81\cdot1024\cdot81.
\]
Поскольку
\[
81=9^2,\qquad 1024=32^2,\qquad 81=9^2,
\]
получаем
\[
\sqrt{9^2\cdot32^2\cdot9^2}=9\cdot32\cdot9=2592.
\]
Результат неотрицателен, как и должен быть квадратный корень.

\textcolor{teal!60!black}{Итог:}
\[
\boxed{2592}.
\]

\textbf{Задача на закрепление.}

Найдите значение
\[
\sqrt{49p^{10}q^4}
\]
при $p=-3$, $q=2$.

\newpage
\phantomsection
\label{task:5}
\section*{Задание №5}

\textbf{Текст задания.}

Упростите выражение
\[
\frac{(a^7b^3)^2}{a^9b^6}.
\]

\textbf{Ваше решение.}

В работе квадрат произведения раскрыт до степеней $a$ и $b$, затем одинаковые множители сокращаются. После сокращения записан результат
\[
0.
\]

\textbf{Ваш ответ.}

\[
0.
\]

\textcolor{red}{Ошибка:} после сокращения потеряна оставшаяся степень $a^5$; сокращение множителей не превращает ненулевую дробь в ноль.

\textbf{Где именно возникает ошибка.}

Корректная часть решения --- раскрытие квадрата:
\[
(a^7b^3)^2=a^{14}b^6.
\]
Следовательно,
\[
\frac{(a^7b^3)^2}{a^9b^6}
=
\frac{a^{14}b^6}{a^9b^6}.
\]

После этого можно сократить $b^6$ и уменьшить показатель степени $a$ на $9$. Первый неверный результат появляется тогда, когда после сокращения записывается $0$ вместо оставшегося множителя $a^5$.

\textbf{Почему это неверно.}

Сократить одинаковый ненулевой множитель в числителе и знаменателе означает разделить и числитель, и знаменатель на этот множитель. Например,
\[
\frac{b^6}{b^6}=1,
\]
а не $0$.

Для степеней с одинаковым основанием действует правило
\[
\frac{a^m}{a^n}=a^{m-n},\qquad a\ne0.
\]
Поэтому
\[
\frac{a^{14}}{a^9}=a^{14-9}=a^5.
\]

Исходный знаменатель содержит $a^9b^6$, поэтому исходное выражение определено только при
\[
a\ne0,\qquad b\ne0.
\]

\textcolor{blue!60!black}{Ключевая идея:} при сокращении одинаковые множители дают единицу, а при делении степеней с одинаковым основанием показатели вычитаются.

\textbf{Исправление от последнего правильного шага.}

Продолжим с дроби
\[
\frac{a^{14}b^6}{a^9b^6}.
\]
Сначала сокращаем $b^6$:
\[
\frac{a^{14}b^6}{a^9b^6}
=
\frac{a^{14}}{a^9}.
\]
Затем вычитаем показатели:
\[
\frac{a^{14}}{a^9}=a^{14-9}=a^5.
\]

\textbf{Полное правильное решение.}

При $a\ne0$ и $b\ne0$:
\[
\frac{(a^7b^3)^2}{a^9b^6}
=
\frac{a^{14}b^6}{a^9b^6}
=
a^{14-9}b^{6-6}
=
a^5b^0
=
a^5.
\]
Здесь $b^0=1$, поэтому множитель $b$ после сокращения исчезает, но это не означает, что всё выражение становится нулём.

\textbf{Проверка.}

Возьмём допустимые значения $a=2$, $b=1$. Исходное выражение:
\[
\frac{(2^7\cdot1^3)^2}{2^9\cdot1^6}
=
\frac{2^{14}}{2^9}
=
2^5.
\]
Формула $a^5$ при $a=2$ также даёт $2^5$. Проверка совпала.

\textcolor{teal!60!black}{Итог:}
\[
\boxed{a^5},\qquad a\ne0,\ b\ne0.
\]

\textbf{Задача на закрепление.}

Упростите выражение
\[
\frac{(c^6d^2)^2}{c^7d^4},\qquad c\ne0,\ d\ne0.
\]

\newpage
\phantomsection
\label{task:6}
\section*{Задание №6}

\textbf{Текст задания.}

Найдите значение
\[
\sqrt{36x^2-60x+25}
\]
при $x=\dfrac12$.

\textbf{Ваше решение.}

В работе выполнено преобразование
\[
\sqrt{36x^2-60x+25}
=
\sqrt{(6x-5)^2}
=
6x-5.
\]
Подстановка $x=\dfrac12$ и окончательный числовой ответ не записаны.

\textbf{Ваш ответ.}

Не указан.

\textcolor{red}{Ошибка:} из квадратного корня извлечено выражение без модуля.

\textbf{Где именно возникает ошибка.}

Последний корректный шаг:
\[
\sqrt{36x^2-60x+25}
=
\sqrt{(6x-5)^2}.
\]
Равенство верно, потому что
\[
(6x-5)^2=36x^2-60x+25.
\]

Первый неверный переход:
\[
\sqrt{(6x-5)^2}=6x-5.
\]
В общем случае справа должен стоять модуль.

\textbf{Почему это неверно.}

Для любого действительного выражения $u$ выполняется
\[
\sqrt{u^2}=|u|.
\]
Квадратный корень всегда неотрицателен. Если $u<0$, то запись $\sqrt{u^2}=u$ дала бы отрицательное число и противоречила бы определению квадратного корня.

В этой задаче при $x=\dfrac12$:
\[
6x-5=6\cdot\frac12-5=3-5=-2<0.
\]
Поэтому модуль особенно важен:
\[
|6x-5|=|-2|=2.
\]
Если использовать запись $6x-5$ без модуля, получится $-2$, что невозможно для значения квадратного корня.

\textcolor{blue!60!black}{Ключевая идея:} полный квадрат под корнем превращается в модуль основания квадрата.

\textbf{Исправление от последнего правильного шага.}

\[
\sqrt{(6x-5)^2}=|6x-5|.
\]
При $x=\dfrac12$:
\[
|6\cdot\tfrac12-5|
=
|3-5|
=
|-2|
=
2.
\]

\textbf{Полное правильное решение.}

Заметим, что
\[
36x^2-60x+25=(6x-5)^2.
\]
Следовательно,
\[
\sqrt{36x^2-60x+25}
=
\sqrt{(6x-5)^2}
=
|6x-5|.
\]
Подставляем $x=\dfrac12$:
\[
|6\cdot\tfrac12-5|
=
|3-5|
=
2.
\]

\textbf{Проверка.}

Подставим $x=\dfrac12$ сразу в исходное подкоренное выражение:
\[
36\cdot\frac14-60\cdot\frac12+25
=
9-30+25
=
4.
\]
Тогда
\[
\sqrt4=2.
\]
Ответ подтверждён.

\textcolor{teal!60!black}{Итог:}
\[
\boxed{2}.
\]

\textbf{Задача на закрепление.}

Найдите значение
\[
\sqrt{49z^2-70z+25}
\]
при $z=\dfrac37$.

\newpage
\phantomsection
\label{task:8}
\section*{Задание №8}

\textbf{Текст задания.}

Найдите значение
\[
\frac{a^{-8}}{a^{-13}}
\]
при $a=-2$.

\textbf{Ваше решение.}

В работе записано
\[
\frac{a^{-8}}{a^{-13}}
=
a^5
=
(-2)^5
=
-512.
\]

\textbf{Ваш ответ.}

\[
-512.
\]

\textcolor{red}{Ошибка:} неверно вычислена пятая степень числа $-2$.

\textbf{Где именно возникает ошибка.}

Последний корректный шаг:
\[
\frac{a^{-8}}{a^{-13}}
=
a^{-8-(-13)}
=
a^5.
\]
Также корректна подстановка
\[
a^5=(-2)^5.
\]

Первый неверный шаг:
\[
(-2)^5=-512.
\]

\textbf{Почему это неверно.}

Пятая степень означает произведение пяти одинаковых множителей:
\[
(-2)^5=(-2)(-2)(-2)(-2)(-2).
\]
Четыре первых множителя дают положительное число:
\[
(-2)^4=16,
\]
после чего
\[
16\cdot(-2)=-32.
\]

Число $512$ равно $2^9$, поэтому оно не может получиться при возведении двойки в пятую степень.

Знак отрицательный, потому что степень нечётная.

\textcolor{blue!60!black}{Ключевая идея:} сначала правильно применить правило деления степеней, затем отдельно проверить модуль и знак степени отрицательного числа.

\textbf{Исправление от последнего правильного шага.}

Из уже полученного выражения
\[
a^5
\]
подставляем $a=-2$:
\[
(-2)^5=-32.
\]

\textbf{Полное правильное решение.}

Поскольку отрицательные показатели требуют $a\ne0$, при $a=-2$ выражение определено.

При делении степеней с одинаковым основанием показатели вычитаются:
\[
\frac{a^{-8}}{a^{-13}}
=
a^{-8-(-13)}
=
a^5.
\]
Подставляем $a=-2$:
\[
(-2)^5=-32.
\]

\textbf{Проверка.}

Можно вычислить исходную дробь через положительные показатели:
\[
a^{-8}=\frac1{a^8},
\qquad
a^{-13}=\frac1{a^{13}}.
\]
Тогда
\[
\frac{a^{-8}}{a^{-13}}
=
\frac{1/a^8}{1/a^{13}}
=
\frac{a^{13}}{a^8}
=
a^5.
\]
При $a=-2$ это снова даёт $-32$.

\textcolor{teal!60!black}{Итог:}
\[
\boxed{-32}.
\]

\textbf{Задача на закрепление.}

Найдите значение
\[
\frac{c^{-4}}{c^{-9}}
\]
при $c=-3$.

\newpage
\phantomsection
\label{task:9}
\section*{Задание №9}

\textbf{Текст задания.}

Найдите значение
\[
\frac{\sqrt{9a^7}\,\sqrt{16b^6}}{\sqrt{a^3b^6}}
\]
при $a=5$, $b=-2$.

\textbf{Ваше решение.}

В работе радикалы преобразованы так, что множитель, связанный с $b^6$, записан как $b^3$ и в числителе, и в знаменателе. После сокращения получено
\[
12a^2,
\]
затем
\[
12\cdot25=300.
\]

\textbf{Ваш ответ.}

\[
300.
\]

\textcolor{red}{Ошибка:} при извлечении квадратного корня из $b^6$ потерян модуль. Финальный ответ совпал с правильным только потому, что одинаковая знаковая ошибка появилась одновременно в числителе и знаменателе.

\textbf{Где именно возникает ошибка.}

До извлечения корней исходное выражение переписано корректно. Первый проблемный переход связан с заменой
\[
\sqrt{b^6}
\]
на
\[
b^3.
\]
При $b=-2$ такая замена неверна.

Действительно,
\[
\sqrt{b^6}=\sqrt{(b^3)^2}=|b^3|.
\]
При $b=-2$:
\[
b^3=-8,
\qquad
|b^3|=8.
\]
Поэтому $\sqrt{b^6}=8$, а не $-8$.

\textbf{Почему это неверно.}

Квадратный корень из квадрата всегда равен модулю основания квадрата:
\[
\sqrt{u^2}=|u|.
\]
Это правило нельзя заменять равенством $\sqrt{u^2}=u$ без проверки знака $u$.

В данном выражении множитель $\sqrt{b^6}$ присутствует и в числителе, и в знаменателе. Если в обоих местах ошибочно заменить его на $b^3$, оба знака изменятся одновременно и затем сократятся. Именно поэтому итоговое число $300$ получилось правильным несмотря на неверное промежуточное правило.

Такое совпадение нельзя использовать как подтверждение корректности решения: при похожем выражении, где такой компенсации нет, ошибка сразу изменит ответ.

\textcolor{blue!60!black}{Ключевая идея:} сначала записать $\sqrt{b^6}=|b^3|$, а уже затем сокращать одинаковый ненулевой множитель $|b^3|$.

\textbf{Исправление от последнего правильного шага.}

При $a=5>0$ и $b=-2\ne0$ имеем
\[
\sqrt{9a^7}=3a^3\sqrt a,
\qquad
\sqrt{16b^6}=4|b^3|,
\]
а также
\[
\sqrt{a^3b^6}
=
\sqrt{a^3}\,\sqrt{b^6}
=
a\sqrt a\,|b^3|.
\]
Тогда
\[
\frac{3a^3\sqrt a\cdot4|b^3|}{a\sqrt a\,|b^3|}
=
12a^2.
\]
Подставляем $a=5$:
\[
12\cdot25=300.
\]

\textbf{Полное правильное решение.}

При заданных значениях $a=5>0$ и $b=-2\ne0$ все радикалы определены, а знаменатель не равен нулю.

Преобразуем каждый корень с учётом знаков:
\[
\sqrt{9a^7}
=
3\sqrt{a^6a}
=
3a^3\sqrt a,
\]
поскольку $a>0$.

Далее
\[
\sqrt{16b^6}
=
4\sqrt{(b^3)^2}
=
4|b^3|.
\]
Знаменатель:
\[
\sqrt{a^3b^6}
=
\sqrt{a^2\cdot a\cdot(b^3)^2}
=
a\sqrt a\,|b^3|.
\]
Следовательно,
\[
\frac{\sqrt{9a^7}\,\sqrt{16b^6}}{\sqrt{a^3b^6}}
=
\frac{3a^3\sqrt a\cdot4|b^3|}{a\sqrt a\,|b^3|}
=
12a^2.
\]
Подставляем $a=5$:
\[
12\cdot5^2=12\cdot25=300.
\]

\textbf{Проверка.}

Проверим прямой подстановкой. При $a=5$, $b=-2$:
\[
\sqrt{16b^6}=4\sqrt{(-2)^6}=4\sqrt{64}=32.
\]
Кроме того,
\[
\sqrt{a^3b^6}
=
\sqrt{5^3\cdot64}
=
\sqrt{8000}
=
40\sqrt5,
\]
а
\[
\sqrt{9a^7}
=
3\sqrt{5^7}
=
375\sqrt5.
\]
Тогда
\[
\frac{375\sqrt5\cdot32}{40\sqrt5}
=
\frac{375\cdot32}{40}
=
300.
\]

\textcolor{teal!60!black}{Итог:} числовой ответ $300$ верен, но промежуточное извлечение корня должно выполняться через модуль:
\[
\boxed{300}.
\]

\textbf{Задача на закрепление.}

Найдите значение
\[
\frac{\sqrt{25m^5}\,\sqrt{36n^6}}{\sqrt{m^3n^6}}
\]
при $m=4$, $n=-3$.

\vfill
\begin{center}
\small Лёвин Артём Александрович эксклюзивно для Ярослава Верещагина
\end{center}

% Краткое сообщение ученику:
% Ярослав, основные трудности в этой работе связаны с правилами действий со степенями и с извлечением квадратного корня из квадрата. В заданиях 2, 4, 5, 6, 8 и 9 стоит особенно внимательно различать правила для степени степени, умножения и деления степеней, а при извлечении корня учитывать модуль. В девятом задании итоговое число получилось верным, но промежуточное преобразование с отрицательным значением переменной было некорректным.
% Репозиторий: students/jaroslav_vereschagin/review_docs/review_14.09.26.tex

\end{document}
